\section{Introduction}\label{sec:intro}

\subsection{The problem}

Can a deterministic multitape Turing machine transpose a $d\times d$ matrix with $b(d)=\max(1,\clog d)$-bit entries, written
row by row, in time linear in the input length $L=d^2b(d)$? This is the statement~(T) of the first paper of this
series~\cite{FirstPaper2026}. By a reduction of Harvey and van der Hoeven~\cite[Cor.~6.9]{HvdH2025}, a linear-time
multiplication machine would give a linear-time transposer, so a proof of $\neg$(T) would give the first superlinear lower
bound for off-line integer multiplication on multitape Turing machines, and a proof of~(T) would remove the data-movement
obstacle to linear-time multiplication. The question is open, and this paper does not decide it.

For oblivious machines the fourth paper of the series~\cite{FourthPaper2026} reduced $\neg$(T) to two statements of different
kinds (its Theorem~D, \cref{thm:D4} here). \emph{Depth reduction}: the time--space dependency DAG of every oblivious linear-time
run can be cut to bounded depth by a genie that reveals $\eps L$ bits. At the level of head walks this is the trajectory cut
problem TCP: deleting $\eps T$ edges of the walk DAG of $k$ walks of length~$T$ should leave every path with boundedly many long
hops. \emph{The crux}: index codes for the side information of the remaining bounded-depth cones need more than $\eps L$ bits.
The fourth paper proved TCP for home-based work heads, whose legs all begin or end near fixed homes, for every coupling of the
heads (its Theorem~B), and it identified as the sharpest open instance the \emph{deep-nesting pair}: two copies of a walk nested
at $m+1$ levels, shifted against each other in time, with a conjecture (A$'$) predicting depth-robustness as $m$ grows.
\emph{Directed Recovery Structures in Tape Computations}~\cite{DirectedRecovery2026} showed that the cone relations of every
linear-time schedule are paired-distance separated after a paid trimming, and isolated the crux as one machine-free question
\cite[Question~9.1]{DirectedRecovery2026}: do all storage codes on bounded-degree recovery relations with this separation have
rate at most $\frac12+o(1)$?

This paper is about the depth side, and about how it meets the coding side. Its main theorem replaces homes by
\emph{owners}, the turning points that hops are nested around, and proves TCP for all walks of bounded nesting load.

\subsection{Results}

\paragraph{Bounded nesting (\cref{sec:nesting}).} The \emph{owner} of a long hop of a walk is the first time inside it at which
the walk is farthest from the hop's cell. The \emph{owner lemma} (\cref{lem:owner}): if $\ell$ nested hops of one side have
pairwise distinct owners, the cell of the innermost hop is entered at least $(\ell+1)/2$ times. So walks that enter every cell
at most $V$ times have \emph{owner-load} at most $2V-1$: at every time, each head has at most $2V-1$ distinct owners among its open
hops of each side. \emph{Theorem~F} (\cref{thm:F}): for $k$ walks of length~$T$, coupled arbitrarily, and any set of their long
hops of owner-load at most~$\ell$, some $\eps T$ edges cut the DAG of these hops to depth $2^{O(k^3\ell^3/\eps)}$, independently
of homes, drift and scales. The label is the fourth paper's band label run on virtual heads, one per head, side and nesting
depth, with bands given by the heights of families of hops with a common owner. Consequently TCP holds for walks of bounded
entry (\cref{cor:entry-tcp}), and in every schedule the depth-robust part, if any, lives on fewer than $k\eps T$ cells that are
entered more than $\lceil1/\eps\rceil$ times, while the rest is cut to depth $2^{O(k^3/\eps^4)}$ (\cref{cor:capacity}). The
deep part must be allowed to relay information: for walks without pauses, forbidding it would already imply TCP for one head
fewer (\cref{prop:relays}).

\paragraph{One arbitrary head (\cref{sec:onehead}).} \emph{Lemma~W} (\cref{lem:W}) cuts a single walk with a strip label, anchored
to cells: $\eps T$ edges and depth less than $16/\eps^2+2/\eps$. \emph{Theorem~M} (\cref{thm:M}) combines it with Theorem~F:
one arbitrary walk, together with walks of owner-load at most~$\ell$ whose families have heights in $q$ dyadic windows of
width~$h$, is cut to a depth that depends only on $(k,\ell,q,h,\eps)$, not on the arbitrary walk or on~$T$; with one window the
depth is $\poly(k,\ell,1/\eps)$ times the square of the ratio of the largest to the smallest family height (\cref{cor:M1}). So a
depth-robust family needs two heads of unbounded owner-load, or one together with bounded heads at unboundedly many separated
scales (\cref{cor:where}). TCP for $k$ walks is, up to a factor~$2$ in~$k$, depth reduction for DAGs with a spine and $k$
non-crossing families of arcs (\cref{prop:stacks}); two families are planar (\cref{prop:sr2}), and three is the first open case.

\paragraph{The deep-nesting pair (\cref{sec:dnp}).} A \emph{turn counter}, the number of turns of either head so far modulo~$q$,
doubly serves every level of both heads at $\log_2q=m+O(\log(1/\theta))$ and cost at most
$\frac{11}{32}\theta n_{\mathrm h}/(m+1)$, at every shift (\cref{thm:counter,cor:counter-served}). Hence Conjecture~A$'$ needs
the constant $C\ge1$ (\cref{cor:C1}), and the pair is cut at depth $\lceil11\cdot2^{m+1}/((m+1)\eps)\rceil-1$, against the fourth
paper's $\lceil12/\eps\rceil^{m+1}-1$ (\cref{cor:counter-cut}): the open question is exponential against subexponential depth in
the nesting depth, with base at most~$2$. We prove that every level class of the pair has exactly $n_{\mathrm h}/(2(m+1))$ hops
(\cref{lem:dnp-classes}). For labels that change only at turns, fine classes are \emph{blind}: on average over shifts, a class with
$g$ levels below it sees a turn of the other head on at most a fraction $(g+3+\lceil\log_2(g+2)\rceil)/(m+1)$ of its hops
(\cref{cor:blind}).
Coarse classes are \emph{visible}: for the block labels that serve one head cheaply, a fraction of order $(m-k)/(m+1)$ of the
apex hops of every coarse class of the other head sees a cancellation at full size, an other-head increment of at most
$-Q-1$, on average over shifts (Theorem~V, \cref{thm:V}). Blindness leaves a safe form of the conjecture, A$''$, which still gives exponential depth (\cref{conj:App}).

\paragraph{Consequences for transposers (\cref{sec:transposers}).} \emph{Corollary~F}: the crux alone, at the depth given by
Theorem~F, excludes oblivious linear-time transposers whose work heads enter every cell boundedly often (\cref{cor:F}); such
heads may reverse unboundedly often, and we know no unconditional lower bound for such machines. \emph{Corollary~F$'$}: an affirmative answer to
Question~9.1 of \emph{Directed Recovery Structures} excludes them (\cref{cor:Fprime}), and \emph{Corollary~F$''$} extends this to
adaptive machines whose traces have entropy at most $\eta L$ for a fixed $\eta<\frac12$, provided every trace of probability at
least $2^{-\eta L}$ has bounded entry (\cref{cor:Fsecond}). Theorem~M gives the analogous
statement when one work head is unrestricted (\cref{cor:M-machine}).

\paragraph{Linear machines (\cref{sec:linear}).} For machines with a data path that is linear over a field, linear depth-TCP,
in which depth is replaced by linear depth and long paths may cancel, together with a linear crux excludes oblivious
linear-time transposers (Theorem~$\Dlin$, \cref{thm:Dlin}). The linear crux can be stated exactly as relative rigidity: the
transposition matrix must not be the sum of a matrix of rank below $\rho L$ and a matrix supported on the bounded-depth ball
pattern of the schedule (\cref{lem:rigidity}). Relative rigidity alone excludes no machine.

\paragraph{The crux in special linear settings (\cref{sec:crux}).} For binary linear codes lifted from a fixed code along a
tower of double covers, an exact recursion with a twist pairing extends the lift identity of \emph{Covers and Local
Generation}~\cite{Covers2026} from symmetric to asymmetric checks, so lifting never raises the rate. Along towers that become
bipartite every code has rate at most~$\frac12$ from that level on, and along towers pulled back injectively from chains with
trivial intersection in free pro-$2$ groups the rate tends to a universal rate at most~$\frac12$, by mod-$2$ Lück approximation
in the form of Jaikin-Zapirain and Souza; for the Clebsch code this limit is~$0$, as \emph{Covers and Local Generation} also
shows. We construct such towers that never become bipartite and have girth tending to infinity (\cref{thm:eqcrux,prop:eqtower}).
On triangle-free graphs~$G$ of maximum degree four, over every field, the kernel of a matrix~$M$ with unit diagonal and support
in~$G$ has dimension at most $\lfloor N/2\rfloor$ when the vertices at which $M$ violates the reciprocity condition of the fourth
paper's Theorem~S are at most two or pairwise at distance at least three (Theorem~SH, \cref{thm:SH}), and in general dimension at
most $\nu(G)+\card H-\rho_3(H)$, where $H$ is the set of these vertices (\cref{cor:SH}).

\paragraph{The construction side and open problems (\cref{sec:construction,sec:open}).} We list what a linear-time transposer
would have to do, among other things carry unbounded nesting and a deep part of vanishing capacity, and the problems on which the
oblivious case rests.

\subsection{How to read this paper: the formal proofs}

Six results are formalized in Lean, two of them in part: the owner lemma (\cref{lem:owner}), the structure of families
(\cref{lem:families}), Theorem~F (\cref{thm:F}) and Lemma~W (\cref{lem:W}) completely; the lemma that cuts are labels
(\cref{lem:labels}) in its path form, its consequence for walk DAGs being deduced in the text (the formal statements of
Theorem~F and Lemma~W include it for their labels); and the label arithmetic of the turn counter (\cref{thm:counter}). The
formal Theorem~F holds for any choice of owners strictly inside the hops, of which the owners of \cref{sec:prelim-walks} are
one. For the turn counter, part~(1), the hypotheses
of the Lean theorem on the cover, and the counts of turns and hops, which rest on the fourth paper's Lemma~5.4, are proved in
the text. Everything else is
proved in the text and not formalized; this includes Theorem~M, Corollaries F, F$'$, F$''$ and~M1, the capacity corollary, the stack form of TCP,
the corollaries of the turn counter, the class count, blindness, Theorem~V, Theorem~$\Dlin$, the equivariant crux and Theorem~SH.
Each of them comes with a complete proof, with all constants, and each numbered result states its status as described in
\cref{sec:prelim-lean}. Part~(f) of the equivariant crux follows from a lemma of \emph{Covers and Local Generation}, and a short script in
\leanfile{tools/} checks it independently by one exact rank computation over~$\F_2$. Results of earlier papers that our proofs use are cited with their status there. Lean-verified there: the first paper's
encoding lemma, pass bounds and coding-gap theorems; and the fourth paper's Lemma~Zq, Theorem~S with its matching form, and
rate bounds for bipartite graphs and for the Clebsch graph. With written proofs there: the fourth paper's Theorem~D,
Corollary~B, local recursion, cone lemma, long-edge conversion and Lemma~5.4; Proposition~9.2 of \emph{Directed Recovery
Structures}; and Lemma~6.1 and Theorem~6.2 of \emph{Covers and Local Generation}. Partly formalized there: the paired-distance
theorem and the paid trimming of \emph{Directed Recovery Structures} (its Theorem~3.3 and Corollary~3.8), whose extension to a
frozen adaptive trace is argued in prose (its Remark~3.9). Numerical observations appear only in remarks labelled as numerical,
heuristics are labelled as such, and claims with only a proof sketch appear only in remarks labelled as such; none of them is
used in proofs.

\subsection{Related work}

\paragraph{The series.} The first paper formulated (T), proved in Lean pass and motion bounds, and reduced $\neg$(T) to the Tape
Coding Conjecture~\cite{FirstPaper2026}; the second excluded, in its models, local coding schemes with low-treewidth interaction graphs
and affine limiting motion~\cite{SecondPaper2026}; the third related bit and word transposition~\cite{ThirdPaper2026}; the fourth proved
Theorems D and~B, curved-motion theorems and storage-code bounds~\cite{FourthPaper2026}. \emph{Directed Recovery Structures in
Tape Computations}~\cite{DirectedRecovery2026} proved the paired-distance structure of native cone relations and posed
Question~9.1, and \emph{Covers and Local Generation in Linear Storage Codes}~\cite{Covers2026} studied linear codes along covers.
Memos to the fourth paper proved the half bound for graphs of circular chromatic number below three~\cite{CircularMemo2026},
subcubic bounds beyond scalar recovery~\cite{SubcubicMemo2026}, and the equivalence of TCP with information-cheap depth
reduction~\cite{InfoDepthMemo2026}. This paper is the eighth of the series. Its companion, the seventh, \emph{Nested
Elimination and Half-Rate Bounds for Storage Codes}~\cite{SeventhPaper2026}, published together with it, proves on the coding side that every storage code on a triangle-free graph of maximum degree three,
and on several classes of directed relations without transitive triangles, has rate at most~$\frac12$, that the guessing number
equals the fractional matching number on such graphs, and that the broadcast rate of index coding there is the number of vertices
minus the fractional matching number.

\paragraph{Depth reduction and depth-robust graphs.} The route from small depth to lower bounds goes back to Valiant's depth
reduction for circuits~\cite{Valiant1977}, and the computation-graph method for multitape machines to Hopcroft, Paul and
Valiant~\cite{HPV1977}. Edge-deletion depth reduction was tied to multitape machines by Paul and
Reischuk~\cite{PaulReischuk1980}, and separators of $k$-page graphs, which contain tape computation graphs, by Galil, Kannan and
Szemerédi~\cite{GKS1989separators}, who also constructed $3$-pushdown graphs with nearly linear
separators~\cite{GKS1989pushdown}. Paul, Pippenger, Szemerédi and Trotter separated deterministic from nondeterministic linear
time with segregators of the computation graph~\cite{PPST1983}, and Santhanam related segregators of multi-pushdown graphs to
time and space~\cite{Santhanam2001}; depth reduction at linear cost deletes a fixed fraction of the edges but must leave
constant depth. Depth-robust graphs go back to a question of Stoß studied by Erdős, Graham and Szemerédi~\cite{EGS1975}, and their
pebbling cost was related to depth-robustness by Alwen, Blocki and Pietrzak~\cite{ABP2017}. Expanders with $O(1)$-pushdown and
three-page layouts exist~\cite{DvirWigderson2010,BourgainYehudayoff2013,DSW2016}. We found no work on reducing DAGs of bounded
stack number, or walk DAGs of bounded nesting, to bounded depth by deleting a linear fraction of the edges.

\paragraph{Coding.} The path sums of \cref{sec:linear} are the transfer matrices of algebraic network coding~\cite{KoetterMedard2003};
minrank is that of Haemers and of Bar-Yossef, Birk, Jayram and Kol~\cite{Haemers1979,BBJK2011}; the route from small depth to
index coding goes back to Valiant's depth reduction and Riis's guessing games~\cite{Valiant1977,Riis2007}, and was used in the
language of min-rank by Jukna and Schnitger~\cite{JuknaSchnitger2011}; for static data structures, lower bounds imply
rigidity~\cite{DGW2019}, and for systematic linear data structures the two are equivalent~\cite{NRR2020}; and the relation between
oblivious algorithms and network codes goes back to Adler, Harvey, Jain, Kleinberg and Lehman~\cite{AHJKL2006}. Coding gaps over routing in
undirected networks were found by Braverman and He~\cite{BravermanHe2026} and by Zhang, Li and Li~\cite{ZhangLiLi2026}. The towers of
\cref{sec:crux-towers} build on \emph{Covers and Local Generation}~\cite{Covers2026}, which proves the lift identity for
symmetric binary checks, the monotonicity of normalized nullity under refinements with $p$-group kernels, the half bound for
the inner rank of graph checks, and full inner rank for checks on graphs with a spanning $2$-factor. They also use Cohn's
theory of free ideal rings~\cite{Cohn2006}, the monotonicity argument of Bergeron, Linnell, Lück and
Sauer~\cite{BLLS2014,Kionke2017}, and the mod-$p$ approximation theorems of Jaikin-Zapirain and of Jaikin-Zapirain and
Souza~\cite{JaikinZapirain2024,JaikinSouza2026}, descendants of Lück's approximation theorem~\cite{Luck1994}.

\subsection{Organization}

\Cref{sec:prelim} recalls the model, walk DAGs, the decomposition of the fourth paper and the deep-nesting pair.
\Cref{sec:nesting} proves the owner lemma and Theorem~F, \cref{sec:onehead} Lemma~W and Theorem~M, and \cref{sec:dnp} the results on
the deep-nesting pair. \Cref{sec:transposers} derives the consequences for transposers, \cref{sec:linear} treats linear machines,
and \cref{sec:crux} proves the two cases of the crux. \Cref{sec:construction} collects what a construction would need, and
\cref{sec:open} the open problems. The appendices index the Lean declarations and explain how to build and check the library.
